\section{Background and Related Work}
\label{sec:background}

\subsection{The Searchable-PDF Problem}
The PDF format~\cite{pdf32000} can represent text either as encoded glyphs with
positioning operators or as raster images. Only the former is extractable. A
``scanned'' PDF is typically a sequence of full-page images with no text
operators, so any tool relying on the document's text stream---including the
browser's find feature---sees an empty page. Making such a document searchable
requires \emph{recovering} the text through OCR and re-associating it with the
page.

\subsection{Optical Character Recognition}
OCR converts images of text into character sequences. Tesseract, originally
developed at HP and later open-sourced by Google, is among the most widely used
engines; modern versions use an LSTM-based recognizer and report results as a
hierarchy of blocks, paragraphs, lines, and words, each with a bounding box and a
confidence score~\cite{smith2007tesseract}. \proj{} uses
Tesseract.js~\cite{tesseractjs}, a WebAssembly port that runs entirely in the
browser, which is what makes a fully client-side, offline pipeline possible.

OCR quality is sensitive to input image quality. Low contrast, blur, and low
resolution all degrade recognition, which motivates the preprocessing stage
described in Section~\ref{sec:impl}. Classic techniques such as Otsu's
thresholding~\cite{otsu1979} binarize an image by analyzing its gray-level
histogram; \proj{} performs a related percentile-based contrast stretch before
handing the image to Tesseract's own binarizer.

\subsection{In-Browser PDF Rendering with PDF.js}
PDF.js~\cite{pdfjs} is Mozilla's standards-based PDF renderer written in
JavaScript. It parses a PDF, exposes each page's content, and rasterizes pages
to an HTML5 \code{<canvas>}. Crucially it also provides two capabilities this
project depends on: \code{getTextContent()}, which returns positioned text items
for pages that \emph{do} have a text layer, and \code{getOperatorList()}, which
exposes the low-level drawing operators and lets us detect whether a page paints
raster images. PDF.js also defines a \emph{viewport} abstraction that maps the
PDF's coordinate system (points, bottom-left origin) into device pixels, which is
the first link in our coordinate chain.

\subsection{The Chrome Extension Platform (Manifest V3)}
Manifest~V3~\cite{mv3} is the current Chrome extension architecture. The
components relevant here are:
\begin{itemize}
  \item \textbf{Service worker} --- an event-driven background script (replacing
        persistent background pages) used here to register network rules on
        install.
  \item \textbf{\code{declarativeNetRequest}}~\cite{dnr} --- a declarative API
        for redirecting or blocking requests without intercepting their contents,
        used to reroute PDF navigations to the custom viewer.
  \item \textbf{Content scripts} --- scripts injected into web pages, reserved
        here for future handling of image text on ordinary pages.
  \item \textbf{Web-accessible resources} and a \textbf{content security policy}
        permitting \code{wasm-\allowbreak unsafe-\allowbreak eval}, required so the bundled
        WebAssembly OCR core can execute.
\end{itemize}

\subsection{Client-Side Storage}
Because OCR is computationally expensive, recovered text is cached in
IndexedDB~\cite{indexeddb}, a transactional, asynchronous, origin-scoped
key--value store suited to structured data. Caching is what turns a slow first
read into an instant subsequent one.

\subsection{Positioning Relative to Existing Approaches}
Server-side services (e.g.\ commercial PDF tools and the open-source
\code{ocrmypdf}) produce a new searchable PDF but require an upload and a
round-trip, and they replace the user's in-browser reading flow. Browser-native
PDF viewers render scanned PDFs but cannot search them. \proj{} occupies the gap
between these: it keeps the live, in-browser experience while adding searchability
locally and privately, and it reuses the browser's own search rather than
shipping a new one. The trade-off is that recognition runs on the client's CPU;
the architecture in Section~\ref{sec:design} is largely a response to making that
trade-off acceptable.
